% GENERATED FILE. Edit canonical lessons, then run npm run generate:books.
% canonical-source-hash: fnv1a64:b00ce3f6828d4671
% canonical-lessons: GE-C369-risiko, GE-C369-riskieren, GE-C369-absichern, GE-C369-garantie, GE-C369-garantieren, GE-R369-review-362-369

\chapter{Ein Risiko eingehen}
\label{ch:ge-ein-risiko-eingehen}
\hlchaptermodality{369}

\begin{chapteropening}
\textbf{By the end of this chapter:} \emph{I can talk about risks and guarantees, and how to protect myself.}

Complete the last lesson of chapter 369: I can talk about risks and guarantees, and how to protect myself.
\end{chapteropening}

\section[das Risiko]{das Risiko — the risk}

\label{lesson:GE-C369-risiko}



\begin{quote}
Before the new one: say the German for to do without, then the German for to replace.
\end{quote}

\subsection*{You'll want to know: das Risiko}
\textbf{das Risiko} — \textquotedblleft{}the risk\textquotedblright{}.

\begin{quote}
Das Risiko ist mir zu groß.
\end{quote}

\begin{quote}
\emph{The risk is too great for me.}
\end{quote}

Plural: \emph{die Risiken}. \emph{Ein Risiko eingehen} — to take a risk.

\begin{cousinweb}
\emph{Risiko} was borrowed from Italian \emph{risico}, \emph{rischio}, a word of the merchants; EtymDB traces French \emph{risque} to Old Italian \emph{risco}. Further back the origin is not settled.
\end{cousinweb}

\begin{grammarlens}[title={What you've built}]
One more word for telling, explaining and describing at B1.
\end{grammarlens}

\subsection*{Your turn}
\begin{itemize}
\raggedright
  \item \emph{Say it:} \emph{das Risiko}
  \item \emph{Say it:} \emph{das Risiko}, once more
  \item \emph{Say it:} a sentence of your own with \emph{das Risiko}
  \item \emph{Recall:} say the German for to do without, then the German for to replace, then say \emph{das Risiko} again
\end{itemize}

\begin{tcolorbox}[breakable,colback=teal!4,colframe=teal!35!black,title={Before you move on}]
What does das Risiko mean? (\textquotedblleft{}The risk\textquotedblright{}.) Say it once more.
\end{tcolorbox}
\section[riskieren]{riskieren — to risk}

\label{lesson:GE-C369-riskieren}



\begin{quote}
Before the new one: say the German for to replace, then the German for the risk.
\end{quote}

\subsection*{You'll want to know: riskieren}
\textbf{riskieren} — \textquotedblleft{}to risk\textquotedblright{}.

\begin{quote}
Er hat alles riskiert und seine eigene Firma gegründet.
\end{quote}

\begin{quote}
\emph{He risked everything and started his own company.}
\end{quote}

Perfekt: \emph{hat riskiert}, with no \emph{ge-}.

\begin{cousinweb}
\emph{riskieren} was borrowed from French \emph{risquer} \textquotedblleft{}to risk\textquotedblright{}, from \emph{risque} \textquotedblleft{}risk\textquotedblright{}, which EtymDB traces to Old Italian \emph{risco}.
\end{cousinweb}

\begin{grammarlens}[title={What you've built}]
One more word for telling, explaining and describing at B1.
\end{grammarlens}

\subsection*{Your turn}
\begin{itemize}
\raggedright
  \item \emph{Say it:} \emph{riskieren}
  \item \emph{Say it:} \emph{riskieren}, once more
  \item \emph{Say it:} a sentence of your own with \emph{riskieren}
  \item \emph{Recall:} say the German for to replace, then the German for the risk, then say \emph{riskieren} again
\end{itemize}

\begin{tcolorbox}[breakable,colback=teal!4,colframe=teal!35!black,title={Before you move on}]
What does riskieren mean? (\textquotedblleft{}To risk\textquotedblright{}.) Say it once more.
\end{tcolorbox}
\section[absichern]{absichern — to protect, to safeguard (mostly reflexive, sich absichern)}

\label{lesson:GE-C369-absichern}



\begin{quote}
Before the new one: say the German for the risk, then the German for to risk.
\end{quote}

\subsection*{You'll want to know: absichern}
\textbf{absichern} — \textquotedblleft{}to protect, to safeguard (mostly reflexive, sich absichern)\textquotedblright{}.

\begin{quote}
Bevor ich unterschreibe, will ich mich gut absichern.
\end{quote}

\begin{quote}
\emph{Before I sign, I want to protect myself properly.}
\end{quote}

Separable: \emph{er sichert sich ab}, \emph{er hat sich abgesichert}.

\begin{cousinweb}
\emph{absichern} is built on \emph{sichern} \textquotedblleft{}to make safe\textquotedblright{}, from \emph{sicher} \textquotedblleft{}safe\textquotedblright{}, inherited from Old High German \emph{sihhur}, which goes back to Latin \emph{sēcūrus} \textquotedblleft{}free from care\textquotedblright{} (EtymDB; Lewis \& Short).
\end{cousinweb}

\begin{grammarlens}[title={What you've built}]
One more word for telling, explaining and describing at B1.
\end{grammarlens}

\subsection*{Your turn}
\begin{itemize}
\raggedright
  \item \emph{Say it:} \emph{absichern}
  \item \emph{Say it:} \emph{absichern}, once more
  \item \emph{Say it:} a sentence of your own with \emph{absichern}
  \item \emph{Recall:} say the German for the risk, then the German for to risk, then say \emph{absichern} again
\end{itemize}

\begin{tcolorbox}[breakable,colback=teal!4,colframe=teal!35!black,title={Before you move on}]
What does absichern mean? (\textquotedblleft{}To protect, to safeguard (mostly reflexive, sich absichern)\textquotedblright{}.) Say it once more.
\end{tcolorbox}
\section[die Garantie]{die Garantie — the guarantee, the warranty}

\label{lesson:GE-C369-garantie}



\begin{quote}
Before the new one: say the German for to risk, then the German for to protect.
\end{quote}

\subsection*{You'll want to know: die Garantie}
\textbf{die Garantie} — \textquotedblleft{}the guarantee, the warranty\textquotedblright{}.

\begin{quote}
Auf den Fernseher gibt es zwei Jahre Garantie.
\end{quote}

\begin{quote}
\emph{The television comes with a two-year warranty.}
\end{quote}

Plural: \emph{die Garantien}. \emph{Das ist noch in der Garantie} — that is still under warranty.

\begin{cousinweb}
\emph{Garantie} was borrowed from French \emph{garantie} (EtymDB), from Old French \emph{guarantie} \textquotedblleft{}protection\textquotedblright{}, from \emph{guarantir} \textquotedblleft{}to vouch for\textquotedblright{}, which goes back to a Frankish word; EtymDB links German \emph{gewähren} \textquotedblleft{}to grant\textquotedblright{} to the same family.
\end{cousinweb}

\begin{grammarlens}[title={What you've built}]
One more word for telling, explaining and describing at B1.
\end{grammarlens}

\subsection*{Your turn}
\begin{itemize}
\raggedright
  \item \emph{Say it:} \emph{die Garantie}
  \item \emph{Say it:} \emph{die Garantie}, once more
  \item \emph{Say it:} a sentence of your own with \emph{die Garantie}
  \item \emph{Recall:} say the German for to risk, then the German for to protect, then say \emph{die Garantie} again
\end{itemize}

\begin{tcolorbox}[breakable,colback=teal!4,colframe=teal!35!black,title={Before you move on}]
What does die Garantie mean? (\textquotedblleft{}The guarantee, the warranty\textquotedblright{}.) Say it once more.
\end{tcolorbox}
\section[garantieren]{garantieren — to guarantee}

\label{lesson:GE-C369-garantieren}



\begin{quote}
Before the new one: say the German for to protect, then the German for the guarantee.
\end{quote}

\subsection*{You'll want to know: garantieren}
\textbf{garantieren} — \textquotedblleft{}to guarantee\textquotedblright{}.

\begin{quote}
Ich kann dir garantieren, dass das Essen dort gut ist.
\end{quote}

\begin{quote}
\emph{I can guarantee you that the food there is good.}
\end{quote}

Perfekt: \emph{hat garantiert}. \emph{Garantiert!} — definitely!

\begin{cousinweb}
\emph{garantieren} was borrowed from French \emph{garantir} \textquotedblleft{}to guarantee\textquotedblright{}, from Old French \emph{guarantir} (EtymDB), a verb built on a Frankish word.
\end{cousinweb}

\begin{grammarlens}[title={What you've built}]
One more word for telling, explaining and describing at B1.
\end{grammarlens}

\subsection*{Your turn}
\begin{itemize}
\raggedright
  \item \emph{Say it:} \emph{garantieren}
  \item \emph{Say it:} \emph{garantieren}, once more
  \item \emph{Say it:} a sentence of your own with \emph{garantieren}
  \item \emph{Recall:} say the German for to protect, then the German for the guarantee, then say \emph{garantieren} again
\end{itemize}

\begin{tcolorbox}[breakable,colback=teal!4,colframe=teal!35!black,title={Before you move on}]
What does garantieren mean? (\textquotedblleft{}To guarantee\textquotedblright{}.) Say it once more.
\end{tcolorbox}
\section[(dialogue)]{Words from chapters 362–369 — a review}

\label{lesson:GE-R369-review-362-369}



\begin{quote}
Say the German for voluntary work and unpaid.
\end{quote}

\subsection*{Your turn}
Say these aloud:

\begin{itemize}
\raggedright
  \item voluntary work, unpaid, voluntary, to donate, to get involved
  \item environmental protection, climate change, sustainable, to recycle, to consume
  \item the condition, the precondition, to require, provided that, as far as
  \item unless, except when, otherwise, apart from that, if need be
  \item the hypothesis, in theory, under certain circumstances, if applicable, at least
  \item possibly, maybe, at any rate, after all, perhaps even
  \item instead, instead of, without, to do without, to replace
  \item the risk, to risk, to protect, the guarantee, to guarantee
\end{itemize}

\begin{tcolorbox}[breakable,colback=teal!4,colframe=teal!35!black,title={Before you move on}]
Voluntary work? (\textbf{das Ehrenamt}.) Unpaid? (\textbf{ehrenamtlich}.) Voluntary? (\textbf{freiwillig}.) To donate? (\textbf{spenden}.) To get involved? (\textbf{engagieren}.) Environmental protection? (\textbf{der Umweltschutz}.) Climate change? (\textbf{der Klimawandel}.) Sustainable? (\textbf{nachhaltig}.) To recycle? (\textbf{recyceln}.) To consume? (\textbf{verbrauchen}.) The condition? (\textbf{die Bedingung}.) The precondition? (\textbf{die Voraussetzung}.) To require? (\textbf{voraussetzen}.) Provided that? (\textbf{sofern}.) As far as? (\textbf{soweit}.) Unless? (\textbf{es sei denn}.) Except when? (\textbf{außer wenn}.) Otherwise? (\textbf{andernfalls}.) Apart from that? (\textbf{ansonsten}.) If need be? (\textbf{notfalls}.) The hypothesis? (\textbf{die Hypothese}.) In theory? (\textbf{theoretisch}.) Under certain circumstances? (\textbf{unter Umständen}.) If applicable? (\textbf{gegebenenfalls}.) At least? (\textbf{wenigstens}.) Possibly? (\textbf{eventuell}.) Maybe? (\textbf{möglicherweise}.) At any rate? (\textbf{zumindest}.) After all? (\textbf{immerhin}.) Perhaps even? (\textbf{womöglich}.) Instead? (\textbf{stattdessen}.) Instead of? (\textbf{anstatt}.) Without? (\textbf{ohne dass}.) To do without? (\textbf{verzichten}.) To replace? (\textbf{ersetzen}.) The risk? (\textbf{das Risiko}.) To risk? (\textbf{riskieren}.) To protect? (\textbf{absichern}.) The guarantee? (\textbf{die Garantie}.) To guarantee? (\textbf{garantieren}.)
\end{tcolorbox}
